\documentclass[11pt]{article}
\usepackage[utf8]{inputenc}
\usepackage{amsthm}



% start: P R E A M B L E ------------------------ P R E A M B L E ------------------------ P R E A M B L E ------------------------

%\input{tikzpreamble.tex}

\title{A Complete Probabilistic Semantics for Lambek Calculus}
\author{}
\date{}

\usepackage[backend=biber,style=authoryear,natbib=true]{biblatex}


%\usepackage{/Users/lachlanmcpheat/Dropbox/Tex/myStyle}

\renewcommand{\P}{\mathcal{P}}

\usepackage[utf8]{inputenc}
\usepackage[T1]{fontenc}
\usepackage{textcomp}
\usepackage[english]{babel}
\usepackage{amsmath, amssymb}
%\usepackage[left=1in, right=1in]{geometry}
\usepackage{tikz}
%\usepackage{xfp, siunitx}
\usepackage{pgfplots}
\usepackage{proof}
\usepackage{dsfont}

% support for linguistics examples
\usepackage{gb4e}
\noautomath % gb4e allows for the use of _ and ^ outside of math mode by default. \noautomath disables this feature.

% figure support
\usepackage{import}
\usepackage{parskip}
\usepackage{xifthen}
\pdfminorversion=7
\usepackage{pdfpages}
\usepackage{transparent}

\newtheorem{proposition}{Proposition}
\newtheorem{definition}{Definition}
\newtheorem{theorem}{Theorem}[section]
\newtheorem{lemma}{Lemma}

\addbibresource{references.bib}

\title{A Complete Probabilistic Semantics for Lambek Calculus}
\author{Aadhavan Mamallan}
\date{November 2025}


\begin{document}



\maketitle

\section{Plan}

We extend the capabilities of Lambek Calculus (LC), the logic of syntactic structure used to model natural language. We look into formalising a probabilistic semantics for LC, in the light of probabilistic methods used in natural language processing with great success. We begin by focusing on parsing, the idea of breaking down a sentence into its syntactic structure. The main limitation parsers based on LC is that a sentence may have ambiguities, resulting in several possible meanings for the sentence. With a probabilistic parser, ambiguities can be resolved by picking the derivation of a sentence with the highest probability.

We start by constructing a formal way of assigning probabilities to words and extending this to sentences (these probabilities can be found by training on a large corpus of text). We then prove the soundness and completeness of LC with respect to these semantics. 

Once this is done, we can compare with existing probabilistic semantics. PCFGs (Probabilistic Context Free Grammars) are a notable example. The equivalence between the expressive powers of Lambek calculus and CFGs was proven by Pentus. We can look into proving a similar equivalance between our probabilistic semantics for LC and PCFGs. 

\section{Completeness}

\begin{definition}
A maximally consistent set is a subset $\Gamma$ of $\mathcal{L} \times [0, 1]$ satisfying the following constraints:
\begin{itemize}
    \item if $(A, p), (A, p') \in \Gamma$, then $p = p'$;
    \item for all $(A, p), (B, s) \in \Gamma$,
    $$
    \begin{array}{lcl}
         A \vdash B & \Longrightarrow & p \leqslant s. \\
    \end{array}
    $$
\end{itemize}
\end{definition}

\begin{definition}
  The canonical relation is a ternary relation $R$ on the set of maximally consistent sets defined as follows:
  $$
  R(\Gamma, \Delta, \Theta) := (A, p) \in \Gamma ~\&~ (B, s) \in \Delta ~\&~ (A \circ B, q) \in \Theta \Rightarrow q \leqslant p \cdot s.
  $$
\end{definition}


\begin{theorem}[left application]
$A \circ (A \setminus B) \models
B$
\end{theorem}

\begin{proof}

Let $x \in \Sigma^\ast$. We need to show that
$v(x, A \circ (A \setminus B)) \leqslant v(x, B)$. Note that
$$ 
\begin{array}{rl}
     v(x, A \circ (A \setminus B)) &  = \\
     \max \{ v(y, A) \cdot v(z, A \setminus B) \mid R(y, z, x) \} & = \\
     \max \{ v(y, A) \cdot \min \{ \frac{v(b, B)}{v(a, A)}\mid R(a, z, b),
  \mbox{with $b \in \mathcal{C}$} \} \mid R(y, z, x) \} \\
  \leq \max \{ v(y, A) \cdot \frac{v(x, B)}{v(y, A)} \mid R(y, z, x) \} \\
  =  v(x, B)
\end{array}
$$
If, for fixed $y, z \in \Sigma^\ast$,
$$
\begin{array}{lcl}
\min \{ \frac{v(b, B)}{v(a, A)} \mid R(a, z, b),
  \mbox{with $b \in \mathcal{C}$} \} & = & \frac{v(x, B)}{v(y, A)}
\end{array}
$$
then these $y, z \in \Sigma^\ast$ contribute the value $v(x, B)$ to the
calculation of the maximum.  If, on the other hand, for fixed
$y, z \in \Sigma^\ast$, there exist $a$ and $b$ with $R(a, z, b)$ for which
$$
\frac{v(b, B)}{v(a, A)} < \frac{v(x, B)}{v(y, A)},
$$
then, due to monotonicity of multiplication, these $y, z \in \Sigma^\ast$
contribute a value strictly less than $v(x, B)$ to the calculation of the
maximum.  It follows that the maximum cannot exceed $v(x, B)$, as required.
Checking of the rest of axioms and rules of inference is left to the reader.
\end{proof}


\printbibliography

\end{document}
